\documentclass[../readings.tex]{subfiles}

\begin{document}

\subsection{Jacobi and Gauss-Seidel Methods}
\label{sub:jacobi-gs}

Standard stationary iterations based on diagonal and triangular splittings.

\subsubsection{Jacobi Iteration}

Decomposes $\bA$ into diagonal $\bD$ and remainder $\bR$: $\bA = \bD + \bR$.

\addDef{Jacobi Scheme}{%
\begin{align*}
    \bx^{(k+1)} = \bD^{-1}(\bb - \bR\bx^{(k)})
\end{align*}
\begin{itemize}
    \item \textbf{Component form:}\enspace $x_i^{(k+1)} = \frac{1}{a_{ii}}\left(b_i - \sum_{j \neq i} a_{ij}x_j^{(k)}\right)$.
    \item \textbf{Concurrency:}\enspace All components update independently using values from the \emph{previous} step.
\end{itemize}
}
\label{def:jacobi-iteration}

\addRemark{%
\textbf{(Jacobi performance)}\enspace
Jacobi iteration is \textbf{inherently parallel}, as all components can be updated simultaneously. However, it often exhibits slow convergence and requires double buffering of the solution vector.
}

\subsubsection{Gauss-Seidel Iteration}

Splits $\bA = \bL_* + \bU_*$, where $\bL_*$ includes the diagonal.

\addDef{Gauss-Seidel Scheme}{%
\begin{align*}
    \bL_* \bx^{(k+1)} = \bb - \bU_* \bx^{(k)}
\end{align*}
\begin{itemize}
    \item \textbf{Component form:}\enspace $x_i^{(k+1)} = \frac{1}{a_{ii}}\left(b_i - \sum_{j < i} a_{ij}x_j^{(k+1)} - \sum_{j > i} a_{ij}x_j^{(k)}\right)$.
    \item \textbf{Update rule:}\enspace Use updated values \emph{immediately} within the same sweep.
\end{itemize}
}
\label{def:gs-iteration}

\addRemark{%
\textbf{(Gauss-Seidel performance)}\enspace
Sequential dependencies restrict efficient parallelization, but Gauss-Seidel typically converges \textbf{twice as fast} as Jacobi and allows for in-place updates.
}

\subsubsection{Convergence Conditions}

\addThm{Stationary Iteration Matrices}{%
For $\bA=\bD+\bL+\bU$, where $\bL$ and $\bU$ are the strict lower and upper triangular parts,
\begin{align*}
    \bT_J = -\bD^{-1}(\bL+\bU)
\end{align*}
is the Jacobi iteration matrix, and
\begin{align*}
    \bT_{GS}=-(\bD+\bL)^{-1}\bU
\end{align*}
is the Gauss-Seidel iteration matrix. By Thm~\ref{thm:iteration-convergence}, each method converges for every initial guess exactly when the corresponding spectral radius is less than $1$.
}
\label{thm:jacobi-gs-iteration-matrices}

\addExcr{%
\begin{enumerate}
    \item Rewrite Def~\ref{def:jacobi-iteration} in the form $\bx^{(k+1)}=\bT_J\bx^{(k)}+\bc_J$.
    \item Rewrite Def~\ref{def:gs-iteration} in the form $\bx^{(k+1)}=\bT_{GS}\bx^{(k)}+\bc_{GS}$.
    \item Apply Thm~\ref{thm:iteration-convergence} to prove the convergence criterion in Thm~\ref{thm:jacobi-gs-iteration-matrices}.
    \item Compute $\bT_J$ and $\bT_{GS}$ for $\bA=\begin{pmatrix}2&1\\1&2\end{pmatrix}$.
\end{enumerate}
}
\label{ex:jacobi-gs-iteration-matrices}

\addThm{Diagonal Dominance}{%
If $\bA$ is \textbf{strictly diagonally dominant} ($|a_{ii}| > \sum_{j \neq i} |a_{ij}|$), both Jacobi and Gauss-Seidel converge for any $\bx^{(0)}$.
}
\label{thm:diag-dom-convergence}

\addExcr{%
\begin{enumerate}
    \item For Jacobi, use strict diagonal dominance to show that the infinity-norm row sums of $\bT_J$ are less than $1$.
    \item Conclude $\rho(\bT_J)\leq \|\bT_J\|_\infty<1$.
    \item Repeat the argument for Gauss-Seidel by isolating the largest component of the error equation.
\end{enumerate}
}
\label{ex:diag-dom-convergence}

\addThm{SPD Convergence}{%
If $\bA$ is \textbf{SPD}, Gauss-Seidel is guaranteed to converge. Jacobi may still diverge.
}
\label{thm:spd-gs-convergence}

\addExcr{%
\begin{enumerate}
    \item For $\bA = \begin{pmatrix} 2 & 1 \\ 1 & 2 \end{pmatrix}$, compute one step of Jacobi and Gauss-Seidel from $\bzero$.
    \item Compare $\rho(\bT_J)$ and $\rho(\bT_{GS})$. Does $\rho_{GS} \approx \rho_J^2$ hold?
\end{enumerate}
}

\subsubsection{Successive Over-Relaxation (SOR)}

Accelerates Gauss-Seidel by overshooting the update.

\addDef{SOR Update}{%
\begin{align*}
    x_i^{(k+1)} = (1-\omega)x_i^{(k)} + \omega \left[\text{GS update}\right]
\end{align*}
}
\label{def:sor}

\addThm{SOR Convergence Window}{%
For SPD systems, SOR converges for relaxation parameters
\begin{align*}
    0<\omega<2.
\end{align*}
For model Poisson problems, well-chosen $\omega>1$ can reduce the iteration count substantially compared with Gauss-Seidel.
}
\label{thm:sor-convergence}

\addRemark{%
\textbf{(SOR acceleration)}\enspace
For the 1D Poisson problem on a grid with spacing $h$:
\begin{itemize}
    \item Jacobi and Gauss-Seidel require $O(n^2)$ iterations.
    \item Optimal SOR requires only $O(n)$ iterations.
\end{itemize}
SOR effectively improves the convergence rate from $O(h^2)$ to $O(h)$.
}

\addExcr{%
\begin{enumerate}
    \item \textbf{Optimal Tuning:}\enspace For $\rho_J = 0.95$, calculate the optimal $\omega^* = 2/(1 + \sqrt{1-\rho_J^2})$.
    \item \textbf{Convergence Sweep:}\enspace Implement SOR and plot iterations to convergence vs.\ $\omega \in [1, 2)$. Identify the ``cliff'' at $\omega^*$.
    \item Check numerically that choices $\omega\leq 0$ or $\omega\geq 2$ fail on an SPD Poisson matrix, as predicted by Thm~\ref{thm:sor-convergence}.
\end{enumerate}
}
\label{ex:sor-convergence}

\subsubsection{Stopping Criteria}

\addDef{Iterative Stopping Tests}{%
Common stopping tests include
\begin{enumerate}
    \item absolute residual: $\|\br^{(k)}\| < \varepsilon$,
    \item relative residual: $\|\br^{(k)}\|/\|\bb\| < \varepsilon$,
    \item increment: $\|\bx^{(k+1)}-\bx^{(k)}\|<\varepsilon$.
\end{enumerate}
Residual tests are tied to the equation; increment tests are cheap but can stop prematurely if convergence is slow.
}
\label{def:iterative-stopping-tests}

\addExcr{%
\begin{enumerate}
    \item Why is a small residual not enough to guarantee a small error $\be^{(k)}$? Use Thm~\ref{thm:residual-error-bound}.
    \item Implement a robust solver that combines relative residual with a \texttt{max\_iter} budget.
\end{enumerate}
}
\label{ex:jacobi-stopping}

\end{document}
